\section*{第 \ref{ch:b1:point-kinematics} 章——质点运动学}

\begin{solution}{exo:b1:point-kinematics:1}
$v = 4.0t$：\qty{3.0}{s} 时为 \qty{12}{m/s}；$a = \qty{4.0}{m/s^2}$；
$x(5) = \qty{50}{m}$；$\qty{100}{km/h} = \qty{27.8}{m/s}$，此时 $t = 27.8/4.0
= \qty{6.9}{s}$。
\end{solution}

\begin{solution}{exo:b1:point-kinematics:2}
$\dot r = 0$，$\dot\theta = \qty{3.0}{rad/s}$：$\vect v = r\dot\theta\,\vect e_\theta
= 6.0\,\vect e_\theta$（\unit{m/s}）；$\vect a = -r\dot\theta^2\vect e_r =
-18\,\vect e_r$（\unit{m/s^2}）。匀速\omterm{cor:b1:point-kinematics:circular}{圆周运动}：速率不变，
\omterm{cor:b1:point-kinematics:circular}{向心加速度}。
\end{solution}

\begin{solution}{exo:b1:point-kinematics:3}
$v = 2\pi R/T = 2\pi \times 6.37\times10^6/86164 = \qty{465}{m/s}$；
$a = v^2/R = \qty{0.034}{m/s^2}$，即 $g/290$。
\end{solution}

\begin{solution}{exo:b1:point-kinematics:4}
$v_0\sin\alpha = \qty{7.71}{m/s}$，$v_0\cos\alpha = \qty{9.19}{m/s}$。
最高点位于 $t = 7.71/9.81 = \qty{0.79}{s}$ 时，高度为 $7.71^2/(2 \times 9.81) =
\qty{3.0}{m}$。射程为 $2 \times 9.19 \times 0.786 = \qty{14.5}{m}$。
落地时\omterm{def:b1:point-kinematics:velocity}{速度}为 $(9.19, -7.71)$，大小为 \qty{12}{m/s}，与
水平方向成 $40^\circ$ 角，斜向下。
\end{solution}

\begin{solution}{exo:b1:point-kinematics:5}
$a_T = \qty{1.0}{m/s^2}$，$a_N = 10^2/25 = \qty{4.0}{m/s^2}$，$a =
\qty{4.1}{m/s^2}$。\qty{5}{s} 后：$v = \qty{15}{m/s}$，$a_N = 225/25 =
\qty{9.0}{m/s^2}$，$a = \qty{9.1}{m/s^2}$。
\end{solution}

\begin{solution}{exo:b1:point-kinematics:6}
$r = b\omega t$，$\dot r = b\omega$，$\ddot r = 0$，$\dot\theta = \omega$：
$\vect v = b\omega\,\vect e_r + b\omega^2 t\,\vect e_\theta$，$v = b\omega
\sqrt{1 + \omega^2t^2}$；$\vect a = -b\omega^3 t\,\vect e_r + 2b\omega^2\,
\vect e_\theta$。$2\dot r\dot\theta$ 项含两个相等的部分：
径向\omterm{def:b1:point-kinematics:velocity}{速度}方向的转动速率为 $\dot\theta$（给出
$\dot r\dot\theta\,\vect e_\theta$），向外运动使横向速率 $r\dot\theta$ 增大（再给出
另一个 $\dot r\dot\theta$）。
\end{solution}

\begin{solution}{exo:b1:point-kinematics:7}
$v = \sqrt{(2.0 \times 1.0)^2 + 0.5^2} = \qty{2.1}{m/s}$；$a = R\omega^2 =
\qty{2.0}{m/s^2}$，指向轴线；螺距 $v_z \times 2\pi/\omega =
\qty{3.1}{m}$；$\rho = v^2/a = 4.25/2.0 = \qty{2.1}{m} > R$。
\end{solution}

\begin{solution}{exo:b1:point-kinematics:8}
$v_{\max} = A\omega = \qty{1.0}{m/s}$；$a_{\max} = A\omega^2 = \qty{10}{m/s^2}$；
$v$ 比 $x$ 超前四分之一周期，$a = -\omega^2x$ 与 $x$ 反向。
$\abs v$ 在 $x = 0$ 处最大，$\abs a$ 在端点 $x = \pm A$ 处最大。
\end{solution}

\begin{solution}{exo:b1:point-kinematics:9}
$\dot{\vect e}_r = \dot\theta(-\sin\theta\,\vect e_x + \cos\theta\,\vect e_y)
= \dot\theta\,\vect e_\theta$；$\dot{\vect e}_\theta = \dot\theta(-\cos\theta\,
\vect e_x - \sin\theta\,\vect e_y) = -\dot\theta\,\vect e_r$。于是
$\vect v = \dot r\vect e_r + r\dot\theta\vect e_\theta$，对其再求一次导数，并用乘积法则，
即得\cref{thm:b1:point-kinematics:cylindrical}中的\omterm{def:b1:point-kinematics:velocity}{加速度}。
\cref{thm:b1:point-kinematics:cylindrical}.
\end{solution}

\begin{solution}{exo:b1:point-kinematics:10}
最高点：$v = v_0\cos\alpha$，$a_N = g$，$\rho = v_0^2\cos^2\alpha/g$。
抛出时：$v = v_0$，$a_N = g\cos\alpha$，$\rho = v_0^2/(g\cos\alpha)$。
最高点处的\omterm{def:b1:point-kinematics:frenet}{曲率半径}是抛出点的 $\cos^3\alpha$ 倍：抛物线在
速率最小、且 $\vect g$ 全部为法向的地方弯曲最厉害。
\end{solution}

\begin{solution}{exo:b1:point-kinematics:11}
$\dot x = v - v\cos(vt/R)$，$\dot y = v\sin(vt/R)$；$\ddot x = (v^2/R)
\sin(vt/R)$，$\ddot y = (v^2/R)\cos(vt/R)$。在地面处（$vt/R = 2\pi n$）：
$\dot x = \dot y = 0$，$\vect a = (0, v^2/R)$，方向向上。在顶部
（$vt/R = \pi$）：$\dot x = 2v$，$\dot y = 0$，速率为 $2v$。
\end{solution}

\begin{solution}{exo:b1:point-kinematics:12}
$\ddot r = 0$，$\dot r = u$，$\dot\theta = \omega$：$\vect a = -ut\omega^2\,
\vect e_r + 2u\omega\,\vect e_\theta$。当 $t = \qty{2.0}{s}$ 时：$-2.0\,\vect e_r
+ 2.0\,\vect e_\theta$（\unit{m/s^2}）。第一项是转动的向心项；
第二项 $2\dot r\dot\theta$ 只因
径向方向本身在转动——在地面上，沿直线以恒定速率行走时，
根本没有\omterm{def:b1:point-kinematics:velocity}{加速度}。
\end{solution}

\begin{solution}{pb:b1:point-kinematics:1}
\textbf{1.} $\vect v = R\dot\theta\,\vect e_\theta$；$\vect a = -R\dot\theta^2
\,\vect e_r + R\ddot\theta\,\vect e_\theta$。

\textbf{2.} $\vect T = \vect e_\theta$（$\theta$ 增大的运动方向），
$\vect N = -\vect e_r$（指向 $C$），$\rho = R$。

\textbf{3.} 最高点：$v^2 = 625 - 392 = 233$，$v = \qty{15.3}{m/s}$。侧边处：
$625 - 196 = 429$，$v = \qty{20.7}{m/s}$。

\textbf{4.} $2v\dot v = -2gR\sin\theta\,\dot\theta$，且 $v = R\dot\theta$：
$\dot v = -g\sin\theta$。当 $0 < \theta < \pi$ 时为负（上升途中减速），
下降途中为正。

\textbf{5.} $a_N = v^2/R = v_0^2/R - 2g(1 - \cos\theta)$；最低点
\qty{62.5}{m/s^2}；侧边 $62.5 - 19.6 = \qty{42.9}{m/s^2}$；最高点 $62.5 -
39.2 = \qty{23.3}{m/s^2}$。

\textbf{6.} 最低点：\qty{62.5}{m/s^2}（无切向分量）；侧边：
$\sqrt{42.9^2 + 9.81^2} = \qty{44.0}{m/s^2}$；最高点：\qty{23.3}{m/s^2}，
方向竖直向下、指向 $C$，与竖直方向的夹角为 $0$。

\textbf{7.} 最低点与最高点处垂直（$a_T = 0$）；
任何处都不平行（处处 $a_N > 0$）。

\textbf{8.} $\dot\theta = v/R = \sqrt{v_0^2 - 2gR(1 - \cos\theta)}/R$：
最低点 \qty{2.50}{rad/s}，最高点 \qty{1.53}{rad/s}。

\textbf{9.} $\ddot\theta = \dot v/R = -(g/R)\sin\theta$：这是长为 $R$ 的
单摆方程；轨道上的汽车就是
一个越过顶点的摆锤。

\textbf{10.} $2\pi\sqrt{R/g} = 2\pi\sqrt{10/9.81} = \qty{6.3}{s}$。

\textbf{11.} $25$、$20.7$、$15.3$、$20.7$ 的平均值：约 \qty{20}{m/s}；
周长 $2\pi R = \qty{63}{m}$：约 \qty{3}{s}。

\textbf{12.} $\sqrt{gR} = \qty{9.9}{m/s} < 15.3$：是。最小值：$v_0^2 \geq
gR + 4gR = 5gR = 490$，$v_0 \geq \qty{22.1}{m/s}$。

\textbf{13.} 最低点：$\vect a = 62.5$ 向上，$\vect g$ 向下：
$\abs{\vect a - \vect g} = 62.5 + 9.8 = \qty{72.3}{m/s^2} = 7.4g$。

\textbf{14.} 最高点：两者都向下，$23.3 - 9.8 = \qty{13.5}{m/s^2} = 1.4g$，
把乘客压向座椅（指向圆心，
对倒置的乘客而言即向上）。

\textbf{15.} 侧边：$\vect a = 42.9$ 水平指向 $C$，再加上 $9.8$
向下（切向）；$\vect a - \vect g$ 为水平方向的 $42.9$：$4.4g$
侧向，指向圆心。

\textbf{16.} 失效：最低点 $7.4g$（最高点与侧边均无问题）。

\textbf{17.} 最高点视重为零：$v_{\text{顶}}^2 = gR$；
于是 $v_0^2 = gR + 4gR = 5gR$，最低点 $a_N = 5g$，载荷为
$5g + g = 6g$，与 $R$ 无关。

\textbf{18.} $a_N = v^2/\rho$：$v$ 由高度（能量）决定，$\rho$
由轨道的几何形状决定；同一高度下改变 $\rho$，
只改变 $a_N$。

\textbf{19.} $a_N + g = 3g$：$a_N = 2g$，$\rho = v_0^2/2g = 625/19.6 =
\qty{32}{m}$。

\textbf{20.} $v_{\text{顶}} = \qty{15.3}{m/s}$；$a_N - g = 1.5g$：$a_N =
\qty{24.5}{m/s^2}$，$\rho = 233/24.5 = \qty{9.5}{m}$。

\textbf{21.} $a_N = v^2/\rho = v^2s/A^2 = v^3t/A^2$：它从零开始线性增长，
\omterm{def:b1:point-kinematics:velocity}{加速度}的变化率为常数，
而不是在进入圆弧处从 $0$ 突跳到 $v^2/R$；
颈部会将这一突跳感受为一次撞击。

\textbf{22.} 水滴形：底部汽车\omterm{def:b1:point-kinematics:velocity}{速度}快，轨道宽阔、曲率平缓；
顶部汽车\omterm{def:b1:point-kinematics:velocity}{速度}慢，轨道收紧。全程载荷保持在
$3g$ 附近，并逐渐施加。

\textbf{23.} 视线：水平距离 \qty{50}{m}，高度
\qty{20}{m}：$r = \qty{53.9}{m}$，仰角 $\beta = 21.8^\circ$。水平\omterm{def:b1:point-kinematics:velocity}{速度}的
横向分量为 $v\sin\beta = 15 \times
0.371 = \qty{5.6}{m/s}$；$\dot\theta = 5.6/53.9 = \qty{0.10}{rad/s}$。

\textbf{24.} 从偏离圆心的点观察，视线的角速率
并不恒定（汽车靠近时快，远离时慢）：
角速率恒定的\omterm{def:b1:optical-instruments:camera}{照相机}会逐渐偏离汽车。

\textbf{25.} 在 $C$ 处用\omterm{def:b1:point-kinematics:cylindrical}{极坐标}描述圆周上的\omterm{def:b1:point-kinematics:velocity}{速度}与\omterm{def:b1:point-kinematics:velocity}{加速度}；
用弗勒内分解区分“加速”与“转弯”，以及缓和曲线；
在\omterm{def:b1:optical-instruments:camera}{照相机}处用\omterm{def:b1:point-kinematics:cylindrical}{极坐标}计算跟踪角速率；
速率规律则由能量给出（后续章节）。
\end{solution}
